\BookPart{Computing with molecules}{foundations}{%
What does it mean for a tube of DNA to \emph{compute}? Part~I answers with one
famous experiment and the questions it raised. Adleman's 1994 solution of a
small route-finding problem turned a graph into strands, a search into
chemistry and an answer into a band on a gel. We reconstruct that experiment
step by step, separate the logic of \emph{generate, filter and verify} from the
molecules that implement it, and then count what molecular parallelism really
costs: copies, volume, time and the effort of reading out an answer.

\textbf{Chapters in this part}
\begin{itemize}\setlength{\itemsep}{2pt}
\item \textbf{1\quad Computing with DNA.} How can molecules represent candidates, remove wrong ones and leave a readable answer?
\item \textbf{2\quad Adleman's experiment.} Which strands, enzymes and separations did the founding experiment use, and what did its evidence establish?
\item \textbf{3\quad Combinatorial search and complexity.} Why does an exponential search stay exponential when it is moved into a test tube?
\item \textbf{4\quad Molecular parallelism and resource accounting.} How many molecules, how much volume and how much time does a molecular search need?
\end{itemize}}
